\scanpage{58}
% manuscript line 1
\uline{Corollaire 1} Pour que $u\in L(E,F)$ soit op. de H.-S., il faut et il
suffit que \deleted{$uv$} $vu$ soit op. de Fredholm pour \deleted{tout} $v\in L^2(F,E)$.
\deleted{Alors} \inserted{Alors}

% manuscript line 3
\[
\|u\|_2=\sup_{\|v\|_2\leq1}|\operatorname{Tr}vu|\qquad(v\in L^2(F,E)) \tag*{(26)}
\]

% manuscript line 4
En effet, on a vu que la condition est nécessaire, et
d'autre part la formule (25) s'écrit

% manuscript line 6
\[
\|u\|_2=\sup_{\|w\|_2\leq1}|\operatorname{Tr}w^*u|\qquad(w\in L^2(E,F))
\]

% manuscript line 7
ce qui résulte immédiatement de $\operatorname{Tr}w^*u=(u,w)$. Réciproquement,
si $vu\in L^1(E,E)$ pour \deleted{tout} $v\in L^2(F,E)$,
l'application $v\to vu$ de $L^2(F,E)$ dans $L^1(E)\deleted{(E)}$
est continue en vertu du th. du graphe fermé,
(car elle est continue quand $L^1(E)\deleted{(E)}$ est muni de
la topologie induite par $L(E)$, qui est séparée et moins fine).
Donc la forme linéaire $\bar w\to\operatorname{Tr}w^*u$ \deleted{(\unsure{sur})} \inserted{sur} $\overline{\bar E\overset{(2)}{\otimes}F}$ est
continue, donc provient d'un élément $u_0$ de $\bar E\overset{(2)}{\otimes}F$
par la formule \deleted{\unsure{$\bar w\to(u,$}} $\operatorname{Tr}(w^*u)=(u_0,w)$ \,; faisant
$w=\bar x\otimes y$, on trouve que $(ux,y)=(u_0x,y)$ pour
\deleted{tout} $x\in E,y\in F$, d'où $u=u_0$\unsure{.}, d'où la conclusion.

% manuscript line 8a
% margin note written sideways
\margin{supprimer \unsure{ou mettre} comme \unsure{simple} application }

% manuscript line 18
\uline{Corollaire 2} Pour que $u\in L(E)\deleted{(E)}$ soit op. de Fredholm,
il faut et il suffit que $u=vw$, où $v,w\in L^2(E)$, on
peut choisir $\|u\|_2=\|v\|_2=\sqrt{\|u\|_1}$.

% manuscript line 21
En effet, la suffisance \deleted{\unsure{est contenue}} \inserted{résulte} du th. 4, la
nécessité se voit en écrivant $u=Uh=Uw^2$, où
$w=\sqrt h$ \,; on a manif\unsure{estement} $\|w\|_2^2=\|h\|_1=\|u\|_1$,
et posant $v=Uw$, on aura $\|v\|_2\leq\|w\|_2$, et
\deleted{$\|u\|_1$} $u=vw$. On a d'ailleurs $\|v\|_2=\|w\|_2$ \inserted{$(=\sqrt{\|u\|_1})$}, \deleted{\unsure{sinon}}
on a\deleted{\unsure{urait}} $\|u\|_1\leq\|v\|_2\|w\|_2<\|u\|_1$.

% manuscript line 27
% archive sheet number: 59
